%%% Local Variables:
%%% mode: latex
%%% TeX-master: "../WeizhouPan_cn.tex"
%%% End:

\section{\hei 教育背景}

\tlcventry{2014}{2011}{\href{http://www.scnu.edu.cn}{华南师范大
    学}~\href{http://portal.scnu.edu.cn/wps/portal/cs}{计算机学院}}{研究生（保送）}{计算
  机应用技术}{}{}

\tlcventry{2011}{2007}{\href{http://www.scnu.edu.cn}{华南师范大
    学}~\href{http://portal.scnu.edu.cn/wps/portal/cs}{计算机学院}}{本科生}{网络
  工程}{}{}

\section{\hei 技能}

\cvcomputer{编程语言}{Vibe Coding, Python, C/C++, Java, Shell/Bash, ECMAScript}
           {操作系统}{Linux, Windows, Mac}
\cvcomputer{英语能力}{CET6(513)}
           {开发工具}{Codex, Claude Code, Emacs, Android Studio, Git, Makefile, gdb}
\cvcomputer{文档排版}{\LaTeX{}, Doxygen, Markdown}
           {测试工具}{minunit, nose, doctest}
\cvcomputer{技术栈}{Android, Cocos, Unity, React}
           {人工智能}{Agent 平台、RAG、微调、评测}

\section{\hei 项目经验}

% Center labels and use "Since"
%\tltextstart[base]{\scriptsize}
%\tltextend[base]{\scriptsize}
%\tlsince{Since~}

\subsection{\hei 工作项目}

\tlcventry{2017}{0}{\href{http://tencent.com}{腾讯公司}}{技术专家 / 基础算法一组组长}{}{}
{
  \begin{tightitemize}%
  \item \href{https://qrobot.qq.com/}{小Q机器人二代}、腾讯开心鼠、QQ 小程序、QQ 机器人及
    多个海外项目（M8、M8 Lord、M9、W 等）技术负责人，现任小Q AI 助手算法负责人、
    QQ 技术委员会 AI 场景探索方向负责人及基础算法一组组长，实线管理算法岗 7 人，开发岗 6 人；
  \item QQ 团队 AI Agent 平台 QPilot 技术负责人，负责算法、架构设计、产品与运营，
    带领 5 人团队从 0 到 1 完成第一个版本，使其成为 QQ 的一站式 AI 底座；
  \item 自研 QPilot 系列模型：\textbf{funcall} 工具调用模型（CoT + SFT + DPO），业务评测集单轮
    准确率 87.22\%$\rightarrow$95.49\%、多轮准确率 65.31\%$\rightarrow$90.81\%；\textbf{wiki} 知识库问答方案
    （司内最早的 iWiki RAG 方案），业务数据集通过率 62\%$\rightarrow$98\%，累计知识库问答 48 万+次、
    衍生问答助手 2890 个；\textbf{code} 中文编程模型，通过增量预训练使鸿蒙 ArkTS 补全达到 SOTA，
    AICR 代码评审正反馈率 8.55\%$\rightarrow$49.85\%，覆盖 600+ 项目、61 万次 CR 任务；
  \item 设计「多创建视图 + 统一 API」的 Ops 架构，规避多模型带来的体验破窗；2024 年平台
    API 调用超 2.7 亿次，覆盖司内 845 个部门、2674 个组，公开 Agent 6500+；
  \item 专利 17 项，KM 文章 75 篇，iWiki 文档 428 篇，KM 技术影响力超 99\%；获腾讯卓越研发奖、领域优秀项目奖、技术影响力奖等。
  \item \href{https://www.infoq.cn/article/QXiWs69g24PMk0ldg1to}{2018 年 GMTC 全球大前端会议演讲嘉宾}；
    腾讯 Next 课堂《微信小游戏入门与实战》、《\href{https://ke.qq.com/course/350627?from=800004380}{使用 Cocos Creator 开发微信小游戏 2048}》、《\href{https://ke.qq.com/course/384790?tuin=1b8113f4}{Python 从入门到实战}》等课程主讲老师；腾讯课堂「薪选」课程等级评定专家团成员（客户端通道负责人）。
  \end{tightitemize}}

\tlcventry{2017}{2015}{\href{http://pingan.com}{平安金融壹账通}}{高级开发工程
  师}{}{}
{
  \begin{tightitemize}%
  \item Hyperion跨平台App开发框架核心成员；
  \item 设计开发了React-Native热更新、H5离线缓存框架、三合一更新器等，提升应用动
    态化能力；
  \item 开发团队代码管理工具
    \href{http://hahack.com/work/enterprise-class-git-version-control-1/}{fmanager}
    ，精通团队代码管理，擅长将复杂的代码管理任务\href{http://hahack.com/tags/Git/}{化繁为简}；
  \item 业余时间为团队开发了微信机器人管家、Code Review平台、代码模块可视化
    平台、任意门线上日志分析平台、全自动加班信息统计平台、“百宝箱”内部网盘、
    “星黎殿”部门书签、“壹翎阁”技术论坛等，大幅提升团队效率。
  \item 平安金融壹账通2016年年度优秀员工，平安科技2015年Q3季度杰出贡献奖；
  \item 项目奖 $\times 1$，工具奖 $\times 2$， 创新奖 $\times 4$，申请专利 8 项
    ；
  \item 人工智能兴趣小组负责人，组织多次机器学习相关技术分享；
  \item 担任 3A 技术论坛演讲嘉宾；
  \item 工作期间保持绩效评级 A 等 3 次。
 \end{tightitemize}}

\tlcventry{2015}{2014}{\href{http://baidu.com}{百度国际科技（深圳）有限公司}}{开
  发测试工程师}{}{}
{
  \begin{tightitemize}%
  \item 移动 APP 源码 crash 隐患扫描工具 \href{http://godeyes.duapp.com}{GodEyes-iOS} 开发者（用户包括：百度、阿里、腾讯、京东）；
  \item 开发设计路径规划 BadCase 挖掘平台，挖掘千余条不合理规划路径，提升产品路
    径规划能力；
  \item 开发一系列底图格式质检工具工具、svn监控工具、
    流量分析工具等工具成为了图像组的基础测试工具。
 \end{tightitemize}}

\newpage

\subsection{\hei 科研项目}

\tlcventry{2013}{2009}{\href{http://www.siat.ac.cn/}{中国科学院深圳可视计算与可视分析重点实验室}}{实习生}{}{}
{
  \begin{tightitemize}%
  \item 基于模型驱动和神经网络的材质演变算法；
  \item 基于二维快速傅里叶变换的形状检索；
  \item 基于k-means算法的三维运动数据聚类（代码已\href{https://github.com/wzpan/MotionSegmentation3D}{开源}）；
  \item 为一个街景导航系统开发了二维地图控件和离线地图数据加密模块;
  \item 5篇核心期刊论文。1项发明专利。中国科学院客座学生奖学金。
\end{tightitemize}}

\subsection{\hei 技术社区}

\tlcventry{2021}{2018}{\href{https://www.cocos.com/}{Cocos}}{社区 KOL / Oteam 核心 PMC}{}{}
{
  \begin{tightitemize}%
  \item 聚集司内多个团队协同开发了一套可在多个业务间复用的高性能、高效能 Cocos Creator
    公共组件框架 \href{https://github.com/cocos/cocomat}{cocomat}，并开源贡献给 Cocos 社区；
  \item 向 Cocos 官方提交 PR 修复 JSB 跨层通信内存泄漏（\#1178）；
  \item Cocos 社区 KOL \& 十周年荣誉讲师，2020 年 Cocos 开发者沙龙讲师；
  \item 腾讯 Cocos Oteam 核心 PMC。
  \end{tightitemize}}

\tlcventry{2017}{0}{\href{http://hexo.io}{Hexo}}{维护者}{}{}
{
  \begin{tightitemize}%
  \item 为流行的个人博客引擎 Hexo 开发了一系列渲染器、检索源生成器、Tag 插件和主题。
  \item 2017年受邀正式成为 hexojs 官方 Github 组织成员。
  \end{tightitemize}}

\tlcventry{2014}{2011}{\href{http://www.gzlug.org/}{广州 GNU/Linux用户组}}{活跃成员}{}{}
{
  \begin{tightitemize}%
  \item 举办 2011 年软件自由日（Software Freedom Day）广州地区庆祝活动并做技术分享；
  \item 常年参与线上技术交流活动及线下 Hacking Thursday 活动。
  \end{tightitemize}}

\tlcventry{2014}{2013}{\href{http://bbs.chinatex.org}{ChinaTeX}}{论坛精英}{}{}
{
  \begin{tightitemize}%
  \item 在论坛分享了多篇 \LaTeX{} 的学习心得。
  \item 高校 \LaTeX{} 推广者。开发华南师范大学的硕士生学位 \LaTeX{} 模板
    \href{https://github.com/scnu/scnuthesis}{SCNUThesis}，支持在线编写，并被广东工业大学、首都师范大学等高校引用。
  \end{tightitemize}}

\subsection{\hei 个人开源项目}

\tlcventry{2017}{0}{\href{http://github.com/wzpan/wukong-robot}{wukong-robot}}{作者}{}{}%
  {
\begin{itemize}
\item \textbf{项目简述}：一个简单、灵活、优雅的中文语音对话机器人 / 智能音箱项目，支持 LLM
  多轮对话能力，还可能是首个支持脑机交互的开源智能音箱项目；可运行在 Raspberry Pi 上，开箱即用；
\item \textbf{主要特性}：支持灵活选择离线唤醒引擎、在线 ASR 引擎和 TTS 引擎，支持
  接入聊天机器人，支持无缝联动 HomeAssistant 。支持插件模块化，容易拓展。支持 Muse 脑机头环唤醒。配套后台管理端，开放 API。支持 Docker 安装，还提供了 hass.io 插件。
\item \textbf{项目指数}：7,100+ stars，1,400+ forks，175 watchers。QQ 用户群 900+ 名。技能插件 25 个。
\end{itemize}}

\tlcventry{2014}{2013}{\href{http://github.com/wzpan/qtevm}{QtEVM}}{作者}{}{}%
  {
\begin{itemize}
\item \textbf{项目简述}：欧拉影像放大技术（Eulerian Video Magnification）的开源
  实现；
\item \textbf{主要特性}：首个完整的 C++ 开源实现，能同时放大动作变化和颜色变化；
\item \textbf{项目指数}：178 stars，67 forks。
\end{itemize}}

\tlcventry{2013}{0}{\href{http://github.com/wzpan/hexo-generator-search}{hexo-generator-search}}{作者}{}{}%
  {
\begin{itemize}
\item \textbf{项目简述}：一个能生成 Hexo 站点检索数据的插件；
\item \textbf{主要特性}：同时支持生成 XML 和
  JSON 两种格式，支持定制生成的内容范围；
\item \textbf{项目指数}：494 stars，84 forks。被多个主题内置。NPM 月下载量 2.3 万次。
\end{itemize}}

\section{\hei 主要奖项}

% Restore normal labels
%\tltext{\scriptsize}

{

\tlcventry{2017}{0}{优秀员工奖($\times 7$)，绩效突出贡献奖($\times 4$)，卓越研发奖，领域优秀项目奖，技术影响力奖（个人），用户口碑奖，代码飞跃奖，社交线 STAR 奖($\times 2$)，第 33 期 PGM 闪电奖“超级码力奖”，SNG“超级码丽”奖，多个 BG 级 / 部门级精品项目奖}{腾讯}{}{}{}

\tlcventry{2016}{2015}{年度优秀员工，杰出贡献奖，“一狼当先”创新员工奖，双月创新奖，项目奖($\times 1$)，工具奖($\times 2$)，创新奖($\times 4$)}{平安金融壹账通}{}{}{}

\tlcventry{2013}{2011}{国家奖学金，第十二届“挑战杯”全国大学生课外学术科技作品竞赛银奖，优秀毕业生，优秀毕业论文}{国家级 / 校级}{}{}{}

\tlcventry{2012}{2008}{校一等奖学金($\times 2$)，校二等奖学金($\times 2$)，创新奖学金}{校级}{}{}{}

}

